\chapter{The reinforcement-learning harness}
\label{ch:harness}

The agent learns purely from the sparse task reward; there is no reconstruction or
auxiliary loss on the recurrent network (the VAE, when used, is pretrained separately and
frozen). The harness has been held essentially constant across architectures so that
comparisons are clean, and it is itself one of the program's hard-won results: a stable
recipe for training a bootstrapped recurrent actor--critic on a task with a deceptive
$0.5$ accuracy floor.

\section{Actor: PAC (maximum-a-posteriori policy optimisation $+$ behaviour cloning)}
The actor loss replaces the clipped PPO surrogate with the PAC objective. For the discrete
action set the maximum-a-posteriori-policy-optimisation E-step is closed form,
\begin{equation}
q(a\mid s)=\mathrm{softmax}_a\!\big[\log\pi_{\theta'}(a\mid s)+\bar Q_{\theta'}(s,a)/\eta\big],
\end{equation}
where $\theta'$ is a time-lagged target network, $\bar Q$ is the critic's mean over the
value distribution, and $\eta$ is a temperature. The M-step plus a behaviour-cloning term
gives
\[
\mathcal L_{\text{actor}}=(1-\alpha)\Big(-\textstyle\sum_a q(a\mid s)\log\pi_\theta(a\mid s)\Big)
+\alpha\big(-\log\pi_\theta(a_t\mid s)\big),
\]
with $\eta=0.1$, $\alpha=0.1$. The critic's value drives the actor \emph{directly}; there
is no generalised-advantage estimation.

\section{Critic: distributional quantile regression}
The critic is a distributional QR-DQN critic: it predicts a set of value quantiles and is
trained with the quantile-Huber loss against the one-step distributional temporal-
difference target $r_t+\gamma\,\mathbf V_{\text{dist}}(s_{t+1})(1-d_t)$. The published
network used a $15$-bin categorical critic; in the reproduction we use $5$--$15$ quantiles.
The number of ``particles'' is a deliberate knob (Chapter~\ref{ch:harnessvar}).

\section{Prioritised episode replay}
Fresh on-policy episodes are augmented with replay episodes drawn at episode granularity
$\propto$ priority$^{\alpha}$, priority being the mean per-step quantile-Huber error,
with importance-sampling weights annealed over training. Episode-level replay is essential
on this task: the discriminating ``correct press'' events are rare, and replay keeps them
in the gradient long enough to be learned.

\section{Target network: hard-copy versus EMA}
The target network $\theta'$ supplies both the critic's bootstrap target and the actor's
E-step reference policy. Two update rules are available: a periodic hard copy ($\theta'\!
\leftarrow\!\theta$ every $N$ optimiser steps, the form that produced the original working
checkpoints) and an exponential moving average ($\theta'\!\leftarrow\!d\,\theta'+(1-d)
\theta$ each iteration). The EMA form was added during the reproduction; whether it is
load-bearing is one of the open A/B questions of Chapter~\ref{ch:harnessvar}.

\section{Burn-in}
For the first $\sim\!20$ iterations the agent is forced to wait, training only the critic
on full-length episodes so the value function and recurrent representation warm up on clean
trajectories before the policy acts. This is the descendant of the program's earlier
reconstruction-pretrain phase and materially reduces the early collapse to premature
pressing.

\section{What the harness taught us}
Three lessons recur in the history. First, \textbf{always-wait is a stable harbor}: a
bootstrapped actor on this task will retreat to the safe action that never incurs a
false-alarm penalty unless exploration is actively maintained. Second, \textbf{the target
network is the main stabiliser}: a time-lagged reference for both the critic target and
the E-step reference policy is what keeps the bootstrap from chasing its own tail. Third,
\textbf{a degenerate $0.5$ optimum has two faces} (always-wait and always-press-at-deadline)
and which one a run falls into is seed-dependent --- both are collapses, and escaping them
is a perception/exploration problem, not a critic problem (the critic learns the $0.5$
value perfectly well in both).
